% \appendix
\renewcommand{\thechapter}{A} % Set Appendix Chapter Label to B

\chapter{Detailed Processing Units}
\label{app:detailed_processing_units}

\section{ExecuteController Sub-Modules}
\label{app:execute_controller_modules}

The \texttt{ExecuteController} is a complex module that contains several specialized hardware units to manage the core computation and data transformation pipeline. The primary sub-modules include:

\begin{itemize}
    \item \textbf{Im2Col Module}: A specialized SRAM controller for convolution operations that converts image data into column format for efficient convolution. It manages sliding window operations over input data, handles address generation for accessing input feature maps, and provides specialized data access patterns for CNN operations.
    
    \item \textbf{Queue (sram\_read\_signals\_q)}: A pipeline synchronization buffer that buffers control signals for SRAM read operations, provides timing alignment between address generation and data availability, and handles pipeline delays in the Im2Col data path.
    
    \item \textbf{Queue (mesh\_cntl\_signals\_q)}: A mesh control signal buffer that queues control signals fed into the computational mesh. It synchronizes computation control with data arrival timing and provides pipeline buffering for mesh operations.
    
    \item \textbf{MultiHeadedQueue}: A multi-port queue for command buffering that allows multiple commands to be read simultaneously. This provides flexible command scheduling and issue capabilities, and is essential for supporting out-of-order execution patterns.
    
    \item \textbf{TransposePreloadUnroller}: A command transformation unit that unrolls complex operations into simpler primitive operations, handles matrix transpose operations for the weight-stationary dataflow, and manages preload command sequences for optimal data placement.
    
    \item \textbf{MeshWithDelays (Computational Mesh)}: The main systolic array for matrix computations that performs the actual matrix multiplication operations. It implements the systolic array dataflow (input/weight/output stationary), handles multiple simultaneous matrix multiplications, and manages timing and data flow through the PE array.
\end{itemize}

\section{Scratchpad Integrated Processing Units}
\label{app:scratchpad_units}

The Gemmini Scratchpad contains several hardware units that process data on-the-fly as it is moved to or from memory, offloading these tasks from the host CPU. The primary units include:

\begin{itemize}
    \item \textbf{\texttt{VectorScalarMultiplier (vsm/vsm\_1)}}: Applies scaling operations to data during movement, with separate scalers for scratchpad and accumulator data paths. This is essential for handling quantized data types.
    
    \item \textbf{\texttt{PixelRepeater}}: Handles pixel repetition optimizations for first-layer convolution operations, reducing memory bandwidth for specific layer types.
    
    \item \textbf{\texttt{AccumulatorScale unit (acc\_scale\_unit)}}: A critical component that applies scaling and activation functions (e.g., ReLU) to accumulator data before it is written back to memory.
    
    \item \textbf{\texttt{Normalizer unit (acc\_norm\_unit)}}: Performs normalization operations on accumulator data.
    
    \item \textbf{\texttt{ZeroWriter}}: Generates zero data when the \texttt{all\_zeros} flag is set, avoiding unnecessary memory reads and saving power.
\end{itemize}
